\documentclass[11pt]{article}

\usepackage[margin=0.85in]{geometry}
\usepackage{amsmath,amssymb}
\usepackage{booktabs}
\usepackage{enumitem}
\usepackage[table]{xcolor}
\usepackage{array}
\usepackage{fancyhdr}
\usepackage{microtype}
\usepackage[hidelinks]{hyperref}

\definecolor{courseblue}{HTML}{234E70}
\definecolor{lightblue}{HTML}{EAF2F8}
\definecolor{lightgray}{HTML}{F3F4F6}
\definecolor{darkgray}{HTML}{374151}

\hypersetup{colorlinks=true,linkcolor=courseblue,urlcolor=courseblue}
\setlength{\parindent}{0pt}
\setlength{\parskip}{0.55em}
\setlist{itemsep=0.2em,topsep=0.3em,leftmargin=1.6em}
\setlist[enumerate]{label=\arabic*.,leftmargin=1.8em}
\renewcommand{\arraystretch}{1.25}

\pagestyle{fancy}
\fancyhf{}
\cfoot{\thepage}
\renewcommand{\headrulewidth}{0pt}

\newcommand{\keybox}[1]{%
  \begin{center}
  \fcolorbox{courseblue}{lightblue}{%
    \begin{minipage}{0.92\textwidth}
      #1
    \end{minipage}}
  \end{center}}

\newcommand{\checkitem}{\item[$\square$]}

\title{\textcolor{courseblue}{\textbf{Final Modeling and Simulation Project}}\\[0.3em]
\large A Guide to Developing, Testing, and Understanding a Computational Model}
\author{PHYS 441}
\date{}

\begin{document}

\maketitle
\pagestyle{fancy}
\thispagestyle{fancy}

\keybox{\textbf{The central principle:} A simple model that you understand
deeply, verify carefully, and investigate thoughtfully is preferable to an
ambitious model that you treat as a black box.}

\section{Purpose of the project}

The final project is an opportunity to take responsibility for a complete
modeling and simulation problem. You will choose a system, formulate a focused
question, implement a model, test the implementation, investigate its behavior,
and explain what the results do and do not mean.

Your project does \emph{not} need to be brand new, publishable, or
earth-shattering. You may begin with a model from this course, another course,
your research, a textbook, or a published example. Originality can be valuable,
but it is not the primary goal. The goal is intellectual ownership.

By the end of the project, you should be able to say:

\begin{quote}
I understand what this model assumes, how it is implemented, how I tested it,
where I encountered difficulties, what its results mean, and what its
limitations are.
\end{quote}

This is an \textbf{individual project}. Discussion with classmates is welcome,
but every student must formulate, implement, investigate, and explain an
individual project.

\section{Learning objectives}

The project will give you practice in the full modeling workflow. You will:

\begin{itemize}
  \item translate a broad topic into a specific, answerable question;
  \item identify state variables, parameters, governing equations, and modeling
        assumptions;
  \item convert a mathematical or rule-based model into a working computation;
  \item distinguish verification from scientific interpretation;
  \item design checks that can expose incorrect or unreliable results;
  \item conduct a controlled computational investigation rather than a single
        demonstration run;
  \item diagnose and document difficulties encountered during implementation;
  \item communicate results, uncertainty, and limitations clearly; and
  \item produce code and documentation that another person can understand and
        reproduce.
\end{itemize}

\section{Choosing an appropriate project}

Begin with a question, not merely a topic. A topic such as ``epidemics,''
``orbits,'' or ``random walks'' is too broad to guide an investigation. A useful
project question identifies a relationship, comparison, prediction, or behavior
that the model can examine.

\begin{center}
\begin{tabular}{>{\raggedright\arraybackslash}p{0.29\textwidth}
                >{\raggedright\arraybackslash}p{0.61\textwidth}}
\toprule
\rowcolor{lightgray}
\textbf{Too broad} & \textbf{Focused enough to investigate} \\
\midrule
Simulate an epidemic. & How do recovery rate and initial infection fraction
affect peak infection and outbreak duration in an SIR model? \\
Model an orbit. & How does numerical energy drift depend on time step for Euler,
midpoint, and fourth-order Runge--Kutta orbit integrators? \\
Study random walks. & How does the mean first-passage time depend on boundary
distance for unbiased and biased one-dimensional random walks? \\
Simulate diffusion. & How do grid spacing and time step affect the stability and
accuracy of an explicit finite-difference diffusion model? \\
\bottomrule
\end{tabular}
\end{center}

An appropriate project has all of the following characteristics:

\begin{itemize}
  \item The question is narrow enough to address during the remaining course
        time.
  \item The model has identifiable assumptions, variables, and rules or
        equations.
  \item A small but complete version can be implemented before extensions are
        attempted.
  \item At least one meaningful verification check is possible.
  \item The model supports a controlled comparison or parameter study.
  \item The expected output is interpretable, not merely visually impressive.
\end{itemize}

If you already have a suitable problem from research or another course, you may
use it. If you prefer to extend a model from PHYS 441, that is also an excellent
choice. If you are unsure what to choose, ask for an instructor-suggested
project. These routes are equally legitimate and will be evaluated with the
same criteria.

\section{Required elements of every project}

\subsection{A specific question}

State the question your simulation will investigate. Your question should tell
the reader what will be varied, compared, estimated, or explained. A project may
evolve, but it should always have a question that determines what evidence is
relevant.

\subsection{A clearly described model}

Explain the model independently of the code. Identify:

\begin{itemize}
  \item the physical, biological, social, mathematical, or engineered system;
  \item the variables that define its state;
  \item the parameters and their meaning;
  \item the equations, transition rules, or update rules;
  \item the initial and boundary conditions, when applicable;
  \item the assumptions that make the model manageable; and
  \item important real effects that the model omits.
\end{itemize}

A diagram, flowchart, table, or pseudocode description may be helpful. The model
description must be understandable without reading the source code.

\subsection{An implementation you understand}

You are not required to write every numerical operation from scratch. You may
use standard scientific libraries and may consult existing implementations.
However, you must understand every important part of the computation and clearly
identify any code or method adapted from another source.

Using a library does not remove the need to explain the model. For example, if
you use an ODE solver, you must still explain the equations, initial conditions,
integration interval, important solver settings, and how the returned solution
is interpreted. Existing software may support your understanding; it may not
replace it.

\subsection{Verification}

Verification asks:

\begin{quote}
\textbf{Did I implement the model that I intended to implement?}
\end{quote}

Every project must include at least one concrete verification test. Possible
tests include:

\begin{itemize}
  \item comparison with a known analytic solution;
  \item a hand calculation for the first few update steps;
  \item a limiting case whose behavior is known;
  \item conservation of energy, mass, probability, or another invariant;
  \item convergence under smaller time steps or finer spatial grids;
  \item comparison with an independently written simpler implementation;
  \item unit tests for important computational functions; or
  \item a benchmark value from a reliable source.
\end{itemize}

``The plot looks reasonable'' is not, by itself, verification. State the expected
result before presenting the test, show the evidence, and explain whether the
test passed.

\subsection{A controlled computational investigation}

A single simulation run is not a complete investigation. After the baseline
model works, study its behavior through one or more controlled comparisons:

\begin{itemize}
  \item vary one or two physically or mathematically meaningful parameters;
  \item compare multiple initial conditions;
  \item compare numerical methods, resolutions, or step sizes;
  \item repeat a stochastic simulation with multiple random seeds;
  \item compare alternative model assumptions; or
  \item perform a sensitivity or uncertainty study.
\end{itemize}

For at least one comparison, record a prediction before running the calculation.
Then compare the actual outcome with the prediction. Unexpected behavior is not
automatically a failure; it is often the beginning of the most interesting part
of the investigation.

\subsection{A difficulty and diagnosis record}

Real computational work rarely proceeds perfectly from start to finish. Every
project must document at least one substantive difficulty. Examples include a
bug, instability, excessive run time, unexpected parameter sensitivity,
inconsistent units, a result that looked plausible but was wrong, or a model
assumption that failed.

For the selected difficulty, explain:

\begin{enumerate}
  \item What did you observe?
  \item What possible causes did you consider?
  \item What evidence helped identify the cause?
  \item What did you change, if anything?
  \item What did you learn from the process?
\end{enumerate}

You are not being graded on having a trouble-free project. You are being graded
on responding to difficulties thoughtfully and using evidence.

\subsection{Interpretation and limitations}

Separate what the program produced from what you infer. Explain the important
patterns in the output and connect them to the original project question. Identify
at least two limitations, such as simplifying assumptions, uncertain parameters,
finite sampling, numerical error, restricted dimensionality, missing physics,
or a narrow validation range. Explain how those limitations constrain the
conclusions.

\newpage
\section{Scope: core, investigation, and extension}

Plan the project in three layers:

\begin{description}[style=nextline,leftmargin=1.5em,labelindent=0pt]
  \item[Core model -- required]
  The smallest scientifically complete implementation. It should run, produce
  interpretable output, and support at least one verification test.

  \item[Investigation -- required]
  A controlled study that uses the working model to address the project question.

  \item[Extension -- optional]
  An additional feature attempted only after the core and investigation work.
  Examples include another dimension, a more realistic force law, a second
  numerical method, improved performance, or comparison with data.
\end{description}

Define the core conservatively. A completed and carefully investigated core is
more valuable than an unfinished collection of ambitious features. Extensions
do not compensate for missing verification or unclear interpretation.

\section{Project schedule and milestones}

Project work begins during the week of November 9 and continues through final
exam week. The checkpoints below are designed to reveal problems early and to
protect time for verification, investigation, and writing. They are progress
targets: you should seek help promptly if you are falling substantially behind
them.

\subsection*{Week of November 9: Question, scope, and proposal}

Choose a topic, turn it into a focused question, and define a core project that
can be completed even if an optional extension is not. By the end of the week,
submit approximately one page containing:

\begin{enumerate}
  \item the focused project question;
  \item a short description of the modeled system;
  \item state variables, parameters, and governing rules or equations;
  \item the planned computational method;
  \item expected plots, statistics, or other outputs;
  \item at least one proposed verification test;
  \item the minimum viable core model; and
  \item one possible extension if time permits.
\end{enumerate}

The proposal is a planning document, not an irreversible contract. Instructor
feedback may reduce the scope or sharpen the question.

\subsection*{Week of November 16: Model and first implementation}

Translate the project idea into a precise model and begin implementing it. By
the end of the week, be prepared to show:

\begin{itemize}
  \item the governing equations or update rules;
  \item a pseudocode, flowchart, or function-level implementation plan;
  \item one predictable test case or hand-worked update; and
  \item evidence that a meaningful piece of the implementation behaves as
        expected.
\end{itemize}

This is the time to discover whether the proposed model is computationally
manageable. Reduce the scope now if the core calculation is not coming together.

\subsection*{November 23--24: Minimum viable simulation}

Use the class days before the Thanksgiving break to connect the pieces. The core
model should now run from start to finish. Be prepared to show:

\begin{itemize}
  \item one complete baseline run;
  \item at least one diagnostic plot or quantitative summary;
  \item preliminary verification evidence; and
  \item a short list of known problems or uncertainties.
\end{itemize}

Do not wait for a polished result before reporting a problem. The purpose of this
checkpoint is to identify difficulties while there is still time to diagnose
them. The Thanksgiving break is not assumed to be a full project work week.

\subsection*{Week of November 30: Verification and investigation}

Complete the verification tests and the controlled comparisons that address the
project question. Interpret the results as they develop, and keep a record of
failed approaches, numerical difficulties, and model limitations. By the end of
the week, the substantive computation should be nearly complete so that the
remaining time can be used for analysis and communication.

\subsection*{Final exam week: Synthesis and submission}

Finish the analysis, figures, report or notebook, code documentation, and
reproduction instructions. Check that every conclusion is supported by evidence
and that another person could follow and run the work. Add an extension only if
the verified core project and its interpretation are already complete.

\keybox{\textbf{Final deadline: Friday, December 11.} The complete project must
be submitted no later than December 11. This is the final, drop-dead submission
date; plan to finish early enough to allow time for a last reproducibility check.}

\section{Final submission}

Submit the following items:

\begin{enumerate}
  \item \textbf{Project report or computational notebook.} It must include the
        question, model, assumptions, method, verification, investigation,
        results, difficulty record, limitations, and conclusions.
  \item \textbf{Complete source code.} The submitted version must reproduce the
        reported results or clearly explain any external data or setup required.
  \item \textbf{Reproduction instructions.} Include required software, the order
        in which files or notebook cells should be run, and any important run-time
        expectations.
  \item \textbf{Sources and acknowledgments.} Cite data, equations, algorithms,
        documentation, code, and other material that informed the project.
\end{enumerate}

Your final report need not recount every line of code. It should tell a coherent
scientific and computational story: what you asked, what you built, how you know
it works, what you learned, and where confidence in the result ends.

\section{Evaluation rubric}

\begin{center}
\small
\begin{tabular}{>{\raggedright\arraybackslash}p{0.27\textwidth}
                >{\centering\arraybackslash}p{0.10\textwidth}
                >{\raggedright\arraybackslash}p{0.53\textwidth}}
\toprule
\rowcolor{lightgray}
\textbf{Category} & \textbf{Weight} & \textbf{Evidence of strong work} \\
\midrule
Question, model, and assumptions & 20\% & Focused question; clear variables,
parameters, equations or rules; assumptions and omissions identified. \\

Implementation and understanding & 20\% & Coherent implementation; important
steps explained; libraries and adapted material used transparently; student
demonstrates ownership rather than black-box use. \\

Verification and numerical reliability & 20\% & Relevant test with a stated
expected result; evidence presented and interpreted; numerical limitations
addressed where appropriate. \\

Investigation and interpretation & 20\% & Controlled comparisons answer the
project question; results are analyzed rather than merely displayed; claims are
supported by evidence. \\

Difficulties, limitations, and reflection & 10\% & A substantive difficulty is
diagnosed thoughtfully; at least two limitations are connected to the strength
of the conclusions. \\

Communication and reproducibility & 10\% & Clear organization and figures;
complete code and instructions; sources acknowledged; another reader can follow
and reproduce the work. \\
\bottomrule
\end{tabular}
\end{center}

Project novelty, code length, visual complexity, and raw computational cost are
not grading categories. A modest project can earn full credit through depth,
correctness, careful investigation, and clear communication.

\section{Final checklist}

Before submitting, confirm that your project contains all of the following:

\begin{itemize}
  \checkitem A focused question that the simulation can address.
  \checkitem A model description understandable without reading the code.
  \checkitem Clearly stated variables, parameters, assumptions, and omissions.
  \checkitem A complete, runnable core implementation.
  \checkitem At least one concrete verification test.
  \checkitem A controlled parameter, method, initial-condition, resolution, or
             stochastic comparison.
  \checkitem A documented implementation difficulty and evidence-based diagnosis.
  \checkitem Results connected explicitly to the original question.
  \checkitem At least two limitations and their consequences for interpretation.
  \checkitem Complete source code and reproduction instructions.
  \checkitem Citations or acknowledgments for external ideas, data, and code.
\end{itemize}

\keybox{\textbf{Remember:} The purpose of the project is not to prove that your
first idea was correct or that your code worked immediately. The purpose is to
show that you can formulate a model, make it computational, test it, learn from
its failures, and use it responsibly to answer a question.}

\end{document}
